\section*{第 \ref{ch:g9:alternating-current} 章——交流电}

\begin{solution}{exo:g9:alternating-current:1}
\omterm{def:g9:alternating-current:dcac}{直流电}在稳定\omterm{def:g8:voltage-voltmeter:voltage}{电压}下始终朝同一方向流动——
这是\omterm{def:g3:first-electric-circuit:battery}{电池}的工作状态。\omterm{def:g9:alternating-current:dcac}{交流电}在摆动的\omterm{def:g8:voltage-voltmeter:voltage}{电压}下\omterm{def:g9:alternating-current:period}{周期}性反向——
这是插座的工作状态。
\end{solution}

\begin{solution}{exo:g9:alternating-current:2}
$T$：一个完整\omterm{def:g9:alternating-current:period}{周期}的\omterm{def:g2:measuring-time:duration}{时长}；$f$：每秒的\omterm{def:g9:alternating-current:period}{周期}数；
$f = 1/T$ 与 $T = 1/f$。对于 \qty{50}{Hz}：$T =
\qty{0.02}{s}$。当 $T = \qty{0.01}{s}$ 时：$f = \qty{100}{Hz}$。
\end{solution}

\begin{solution}{exo:g9:alternating-current:3}
一百次：每个\omterm{def:g9:alternating-current:period}{周期}内有两次反向，
共五十个\omterm{def:g9:alternating-current:period}{周期}。
\end{solution}

\begin{solution}{exo:g9:alternating-current:4}
$T = 4 \times 5 = \qty{20}{ms} = \qty{0.02}{s}$; $f = 1/0.02
= \qty{50}{Hz}$.
\end{solution}

\begin{solution}{exo:g9:alternating-current:5}
\qty{230}{V} 是\omterm{prop:g9:alternating-current:effective}{有效电压}——即能产生相同\omterm{prop:g8:electric-current-effects:three}{热效应}的稳恒\omterm{def:g9:alternating-current:dcac}{直流电}压值，
也是交流\omterm{def:g8:voltage-voltmeter:voltmeter}{电压表}的读数；
\qty{325}{V} 是峰值，即\omterm{def:g8:voltage-voltmeter:voltage}{电压}摆动实际达到的波峰，
两个方向各每秒出现五十次。
\end{solution}

\begin{solution}{exo:g9:alternating-current:6}
\omterm{prop:g8:electric-current-effects:three}{热效应}与方向无关——电流无论朝哪个方向通过线圈，线圈都会发热，
因此电热水壶可直接使用交变\omterm{def:g8:voltage-voltmeter:voltage}{电压}。
充电是化学反应的逆过程，而化学反应要求
电流方向保持一致，因此充电器先将交变\omterm{def:g8:voltage-voltmeter:voltage}{电压}整为稳定的\omterm{def:g9:alternating-current:dcac}{直流电}压。
\end{solution}

\begin{solution}{exo:g9:alternating-current:7}
\omterm{def:g3:first-electric-circuit:battery}{电池}：\qty{9}{V} 处的一条水平直线。\omterm{def:g9:alternating-current:dcac}{交流电}源：
平滑的波形，在 $+9$ 与
\qty{-9}{V} 之间摆动，每 \qty{20}{ms} 一个\omterm{def:g9:alternating-current:period}{周期}。一眼即可辨别：
平直还是起伏。
\end{solution}

\begin{solution}{exo:g9:alternating-current:8}
旋转的机器自然产生\omterm{def:g9:alternating-current:dcac}{交流电}；只有\omterm{def:g9:alternating-current:dcac}{交流电}能供给
变压器，没有变压器，电能就无法以高\omterm{def:g8:voltage-voltmeter:voltage}{电压}远距离输送，
并以安全的低\omterm{def:g8:voltage-voltmeter:voltage}{电压}进入家庭。
\end{solution}

\begin{solution}{exo:g9:alternating-current:9}
蹬踏较慢时，发电机每秒的少量\omterm{def:g9:alternating-current:period}{周期}使
灯明显闪烁——每个波峰一次闪亮，每个零点一次变暗；蹬得更快时，
脉冲成倍增加，直到人眼无法分辨。
电影的 $24$ 幅画面大致确定了人眼的融合速率：低于
约 \qty{25}{Hz} 的光脉冲时，闪烁令人烦躁；高于此值时，
灯光看起来稳定。
\end{solution}

\begin{solution}{exo:g9:alternating-current:10}
波峰约为 \qty{325}{V}——每秒共一百个，
其中正向、负向各五十个。
\end{solution}

\begin{solution}{exo:g9:alternating-current:11}
$T = 1/60 \approx \qty{0.017}{s}$；峰值 $\approx 120 \times
1.4 \approx \qty{170}{V}$。电热器和\omterm{def:g3:first-electric-circuit:bulb}{灯泡}对此
毫不在意；现代充电器标签上标注“100--240 V，50/60 Hz”，
并能良好地转换——不能适应的是
老式单一\omterm{def:g8:voltage-voltmeter:voltage}{电压}的装置，以及任何
通过计数电网\omterm{def:g9:alternating-current:period}{周期}工作的电动钟：在六十\omterm{def:g8:sound-pitch-loudness:frequency}{赫兹}下会走快。
\end{solution}

\begin{solution}{exo:g9:alternating-current:12}
灯在两个波峰处都变亮，得到 $2 \times 50 = 100$ 个脉冲
每秒。一台 $25$ 帧的相机每
\qty{0.04}{s} 采样一次——每帧包含四个脉冲，但在滚动曝光的各点上，
对齐情况并不相同：相机
的节奏与灯的节奏之间存在微小差异，
于是在相继的各帧中形成缓慢移动的条纹——一种
两个时钟之间的拍。当 \qty{60}{Hz}（$120$ 个脉冲）时，
与 $25$ 帧之间的错配随之改变，条纹的
宽度与\omterm{def:g7:motion-average-speed:formula}{速度}随之改变——图样泄露了当地电网的\omterm{def:g8:sound-pitch-loudness:frequency}{频率}。
\end{solution}

\begin{solution}{pb:g9:alternating-current:1}
\textbf{1.} \omterm{def:g9:alternating-current:dcac}{直流电}，$2.25 \times 2 = \qty{4.5}{V}$：平直线，
对应袖珍\omterm{def:g3:first-electric-circuit:battery}{电池}。
\textbf{2.} 交流；$T = \qty{20}{ms}$，$f = \qty{50}{Hz}$；峰值
$3 \times 2 = \qty{6}{V}$。
\textbf{3.} 直流，\qty{4.5}{V}，\emph{反向}连接——
即正负极对调：对\omterm{met:g9:alternating-current:scope}{示波器}无害，但
颇有教益。
\textbf{4.} 交流；$T = \qty{10}{ms}$，$f = \qty{100}{Hz}$；
峰值 \qty{3}{V}。其每秒一百次的加倍脉冲，
更从容地超出人眼的融合速率，相比 B 的
五十\omterm{def:g9:alternating-current:period}{周期}闪烁。
\textbf{5.} \omterm{def:g8:voltage-voltmeter:voltmeter}{电压表}显示的是有效值——
即与\omterm{prop:g8:electric-current-effects:three}{热效应}等效的直流值；对正弦波而言，它
约为峰值除以 $1.4$：$6 \div 1.4 \approx
\qty{4.3}{V}$。两个数值描述的是同一种摆动。
\textbf{6.} 仅电源 A 可用（方向正确；C 在
重新插接后亦可）。\omterm{def:g9:alternating-current:dcac}{交流电}源与任何\omterm{def:g3:first-electric-circuit:battery}{电池}之间，
都需接入充电器的整流块——
作为变换环节。
\textbf{7.} 市电每秒五十次的摆动接近
心脏自身的电节律，可能扰乱它：同为一\omterm{def:g8:intensity-ammeter:intensity}{安培}，
交变电流比稳恒直流更危险。
\textbf{8.} 正弦曲线——理想的摆动，
由线圈在磁场中以恒定速率转动时自然画出，
即下一章的交流发电机。
\textbf{9.} 每 $4$ 格一个\omterm{def:g9:alternating-current:period}{周期}（\qty{20}{ms}）；波峰
位于上下各 $325 \div 100 = 3.25$ 格处。
\textbf{10.} $f = 1/0.02 = \qty{50}{Hz}$；有效值 $\approx
325 \div 1.4 \approx \qty{230}{V}$——即标签上的数值，
由屏幕上还原而得。
\textbf{11.} \qty{230}{V} 处的一条水平线——表示
等热功率下的稳定性。
\textbf{12.} 例如：“平直线是\omterm{def:g3:first-electric-circuit:battery}{电池}的恒定性——方向单一，
数值不变。波形则是电网的
摆动——方向与数值在不停循环中变化。
而二者之间公平的兑换比率，正是\omterm{prop:g9:alternating-current:effective}{有效
电压}：热效相等，价值便相等。”
\end{solution}
